\documentclass[../main.tex]{subfiles}
\begin{document}
%==========================================TIÊU ĐỀ TẬP========================================
\begin{table}[h]
  \begin{adjustwidth}{-1cm}{}
    \begin{tabular}{>{\centering\arraybackslash}p{6cm}|>{\raggedright\arraybackslash}p{12.5cm}}
      \multirow{5}{6cm}
      % Cột bên trái: ÉPISODE và số tập
      {\\[-40pt]
        \flaregothic\fontsize{20pt}{20pt}\selectfont \centering
        \color{eptype}ÉPISODE\color{black}\\
        \burbank\fontsize{72pt}{72pt}\selectfont
        \color{epnum}125
        \\[18pt]
        % \flaregothic\fontsize{14pt}{14pt}\selectfont
        % \color{epattr}BIENVENUE, \uline{\color{Banpire} Mel} !
        \flaregothic\fontsize{18pt}{18pt}\selectfont
        \color{epattr}ÉPISODE PERDU
        \color{black}
      }
       &
      % Dòng thứ nhất: tên tập tiếng Nhật
      {\fotskip\fontsize{18pt}{18pt}\selectfont \ruby{捕}{つか}まりました。
      }                                           \\[6pt]
      % Dòng thứ hai: tên tập tiếng Anh
       & {\cafeteria\fontsize{18pt}{18pt}\selectfont Jailed for Y*uT**e Crimes}                             \\[4pt]
      % Dòng thứ ba: tên tập tiếng Việt
       & {\vnmsans\fontsize{18pt}{18pt}\selectfont \href{https://youtu.be/jTO5AnmKWCE}{Bạn hù? Bạn vô tù.}} \\[8pt]
      % Dòng thứ tư: Nhân vật xuất hiện
       & {\sen{\fontsize{14pt}{14pt}\selectfont Track used: Cobweb (\ruby{蜘蛛}{くも})}}
      \\[6pt]
      % Dòng cuối cùng: Thời gian diễn ra của tập (thường sẽ là ngày phát hành)
       & {\continuum\fontsize{16pt}{16pt}\selectfont Ngày phát hành: 03/10/2021}                            \\[8pt]
       & (dựa theo bản dịch của \href{https://www.facebook.com/mamisubteam}{MamiSubs - Mami Fansub})        \\[18pt]
    \end{tabular}
  \end{adjustwidth}
\end{table}
%=========================================TRUYỆN CHÍNH========================================
\color{black}

Một buổi sáng đầu tháng Mười, bầu trời thu ở Akiba ảo chuyển thành màu xám nhẹ, không nắng gắt cũng chẳng quá lạnh, cái thời tiết khiến ai cũng muốn lười biếng đôi chút.

Trong văn phòng \hll, Lamy đang ngồi thư thả bên chiếc máy tính của nhân viên, miệng ngân nga vài câu không đầu không cuối: ``Lamy Lamy Lamy\~{} Hôm nay được nghỉ\~{} Xem stream của mình thôi\~{} xem luôn cả của mọi người nào\~{}''

Ngón tay cô lướt chuột nhanh nhẹn, thi thoảng còn bấm ``like'' cho chính kênh mình và vài thước phim của Aki. Tâm trạng nàng Yêu tinh tuyết đang rất tốt.

Bất chợt, từ phía sau, một cô gái tóc vàng, cột thành hai búi hai bên, bất ngờ gào to bằng giọng ngọng líu lưỡi như đang đọc thần chú: ``\textbf{Bông tể hào cựu lược bữa nồi, la rẽ hắt tươi\footnote{Câu gốc của Nene là「タマランチ会長大はしゃぎ！！！！」, một câu vô nghĩa, vần với ``tomaranai'' (không thể chịu đựng nữa). Ở đây người dịch sử dụng một câu vô nghĩa tương đương, dịch ra là: ``Không thể nào chịu được nữa rồi, ta sẽ bắt ngươi!''}!!}''

\img{10-5a}

Lamy hét toáng lên, ghế bật ra phía sau. Nàng chưa kịp hiểu chuyện gì thì đã vội chạy thẳng đi báo với Subaru: ``\textbf{Có người định giết hại em!!}''

Subaru, đang đứng gọt táo bên phòng nghỉ, nghe xong liền phóng thẳng như lực lượng đặc nhiệm, tóm lấy thủ phạm là Nene, kéo thẳng tới một chiếc ``lồng tạm giam'' thứ mà Kuon và Korone đã từng dùng để diễn kịch năm ngoái và vô tình bỏ quên ở văn phòng\footnote{Xem lại {\flaregothic ÉPISODE 40}, quyển số Ba.}.

``\textbf{Thả tôi ra! Tôi chỉ giỡn chơi thôi mà!}'' nàng ``Hải cẩu'' hét lên, hai tay nắm chặt song gỗ, mắt long lanh nước như chim non bị bắt nhốt.

``Mình phải tìm cách vượt ngục thôi\dots\ Đâu thể sống mãi như vầy\dots''

Bỗng từ trong bóng tối của ``buồng giam'', có hai giọng nói cùng vang lên:

``Có vẻ em cũng bị bắt vào đây nhỉ, Nene-chan.''

``Dẹp ý định đó đi.''

Nene giật mình, quay đầu lại. Một người là cậu con trai mà ai cũng biết tiếng vì sự thẳng tính tới mức đôi khi thô lỗ, còn người kia - người bạn đồng thế hệ - là một trong những thành viên ``màu mè'' nhất \hll.

``C-Chúa tể Vượt ngục!?'' nàng ``Hải cẩu'' thốt lên khi thấy Polka.

``Và đồng thời cũng thích màu đỏ,'' Kuon vừa nói vừa chăm chú gõ gì đó vào chiếc laptop đặt trên đùi.

``Ủa, sao anh cũng ở trong này vậy!?'' Nene hỏi, không giấu nổi sự ngạc nhiên.

Cậu Tổng đáp tỉnh bơ, mặt không hiện một chút biểu cảm: ``Anh bảo Subaru-chan cho anh vào đây để dễ tập trung làm việc.''

Polka, đối lập hoàn toàn với vẻ nghiêm túc kia, dang tay phô diễn như nghệ sĩ xiếc đang trong màn trình diễn: ``Chào mừng đến với Gánh Xiếc Tội Phạm!''

``Đừng làm quá nữa\dots'' Kuon lầm bầm, không nề hà gì tới câu nói phóng đại của cô.

Nene ngó ra ngoài, thấy một ``sinh vật ba đầu'' đang gằm gừ canh gác bên ngoài, ánh mắt rực lửa.

``Nhìn kìa! Một con Cerberus đang trông chừng bọn mình kìa!'' nàng Cáo sa mạc lên giọng, đầy khí thế.
``Nó mà tỉa cái là còn cái địa ngục\dots''

``Hay bọn anh còn gọi là còn cái nịt đấy,'' cậu Tổng nói chen vào, ``Với lại\dots''

\img{10-5b}

``\dots Đó chẳng phải Cerberus gì đâu. Chỉ là Watame-chan đang rảnh háng không có gì làm thôi.''

Quả đúng như cậu Tổng nói. Không có con chó ba đầu nào cả, mà chỉ có nàng Cừu - với bộ móng vuốt lắp ở tứ chi - đang bò xuống, gầm gừ như một con chó canh nhà hung ác (mà thực chất là đến việc đó cô cũng chỉ giả vờ làm vậy).

Polka cứng họng, trán lấm tấm mồ hôi. ``Đừng có làm người ta tụt mood như thế chứ Kuon-senpai!''

``Anh chỉ nói sự thật thôi,'' giọng cậu điềm nhiên, không một chút lay động cảm xúc.

Sự tương phản giữa Polka rực rỡ, hoạt ngôn, tưởng tượng như trẻ con và Kuon điềm đạm, thực tế đến mức tàn nhẫn hiện rõ hơn bao giờ hết. Một người đẩy sự hoang tưởng đi xa như truyện tranh, người kia thì kéo lại về thực tại như nhắc: ``Tỉnh mộng đi mấy đứa.''

Nàng Chủ xiếc thở dài thườn thượt trước câu nói phũ như búa bổ của Kuon, không chút cảm xúc, chẳng hề nể nang. Dù vậy, cô vẫn cố lấy lại khí thế, chỉnh lại cổ áo như thể chưa có gì xảy ra, rồi nheo mắt nhìn Nene, giọng kênh kiệu: ``Hạ đẳng như ngươi lo mà chờ mãn hạn tù đi.''

Cậu Tổng cười khẩy, vẫn gõ máy tính như thể đang ghi chép lại buổi họp với Quốc hội Trung ương: ``Làm như thể em có thể vượt ngục dễ dàng ấy. Bao nhiêu lần em thoát nổi chưa? Toàn bị bắt lại.''

Polka khựng lại. Gương mặt đang đắc ý bỗng trùng xuống như vừa bị dí đầu vào gối. Nene nghe vậy, hỏi lại với sự hoài nghi: ``Khoan\dots\ Nếu bà là Chúa tể Vượt ngục, sao vẫn còn ở đây?''

Kuon không ngần ngại liếc sang, nói một câu như đinh đóng cột: `Đó, nghe vô lý chưa?`''

Nàng Cáo xiếc há miệng, không nói được lời nào, ``tắt điện'' toàn tập. Cố chữa cháy, nàng lầm bầm: ``Tại\dots\ tại cai ngục dễ thương quá.''

``Ảo đùa la xi bùa!?'' nàng ``Hải cẩu'' kêu lên, mắt trố ra.

``Nghe vô văn lý vãi chưởng\dots'' cậu lắc đầu. ``Mà em nói cũng đúng, với lại\dots''

*KẸT KẸT*

Âm thanh xe đẩy vang lên từ bên ngoài. Ba người cùng hướng mắt. Nene reo lên: ``Đồ ăn! Đúng rồi! Thời cơ tới!''

``Em định làm gì với chúng?'' Kuon hỏi.

``Dễ lắm! Đánh lạc hướng cai ngục, rồi mình trốn!''

``Không có tác dụng đâu chị bé.'' Polka nhắc nhở nhẹ nhàng như thầy cô giáo dạy Đạo đức.

Kuon chêm vào: ``Thật. Pol-pol thử rồi. Tác dụng bằng mắt.''

``Ể!? Tại sao!?''

Polka đưa tay lên ngực, mắt long lanh: ``Tại bồi bàn dễ thương quá.''

Cậu gật đầu ra chiều đồng cảm: ``Cái này anh phải đồng ý. Nhưng mà\dots\ em dễ siêu lòng thật, Pol-pol ạ.''

``Ý anh là gì chứ!?''

``Bảo sao em không vượt ngục nổi. Chết vì gái là cái chết tê tái mà.''

Vừa lúc ấy, Luna xuất hiện với xe đẩy hai tầng, trên đó là một chai rượu vang mini, vài bát takoyaki còn nóng, một đĩa sashimi và\dots\ một con cá nguyên con quấn khăn ăn ở tầng dưới như chuẩn bị vào tiệc buffet.

\img{10-5c}

``Chị nang đù ăn về cho em nè-nora\~{}!''

Nene liếc sang Polka, rồi chỉ tay về phía Kuon, hạ giọng: ``Kuon-senpai nói đúng quá. Chúa tể Mê Gái còn gì.''

``Đúng chứ?'' Kuon quay sang Luna, người đang đủng đỉnh đầy phấn khởi. ``À, Luna-chan, cho anh xiên bánh takoyaki nha.''

``Có niền-nora\~{}!''

``Em vẫn ngạc nhiên là anh bình tĩnh được đấy\dots'' nàng ``Hải cẩu'' lẩm bẩm.

*EEEE-OOOO* *EEEE-OOOO*

Tiếng còi báo động rú lên.

``\textbf{Gì vậy!?}'' Nene hốt hoảng.

``Báo động có kẻ vượt ngục,'' nàng Cáo xiếc tỉnh bơ nói, khoanh tay lại, như thể cô đã từng trải qua chuyện này.

``\dots Lại em ấy nữa à?'' cậu Tổng thở dài.

``Cá là vượt không thành\dots''

``\dots lần nữa cho mà xem.''

Cảnh bên ngoài nhà giam hiện ra: Rushia bị còng tay bằng lắc bạc, cố vùng vẫy như con cá mắc cạn, vừa la hét: ``\textbf{Á! Thả ta ra!}''

\gif{10-5d}{1}{81}

Nàng ``Hải cẩu'' lo lắng nhìn đàn chị: ``Chị ấy sẽ bị làm gì vậy?''

``Chờ Cerberus - trong ngoặc kép - tra tấn thôi,'' cậu Tổng đáp với giọng điềm nhiên.

``Tra tấn?''

``Nghĩa là sẽ phải nghe cô ta đàm đạo cả đêm,'' nàng Cáo sa mạc giải thích ngắn gọn.

``\dots Trời đất.''

Đúng như lời nàng Chủ xiếc, Watame đang ``tra tấn'' Rushia bằng lời nói: một mạch dài bất tận về cách làm bánh mochi, về thời sự, về nhạc Enka, về stream hôm qua\dots\ giọng đều đều, ánh mắt phấn khích như giáo viên nhận lớp mới. Rushia thì run bần bật, mắt trắng dã, như đang cố sống sót trong cơn mưa bão thông tin.

\img{10-5e}

``'Giá mà em ấy có tai nghe thì đỡ khổ\dots' Kuon chép miệng, rồi tiếp tục đưa mắt về màn hình máy tính.

Trong khi đó, hai nàng thế hệ thứ Năm nhanh chóng tranh thủ cơ hội.

``Cơ hội tới rồi! Chạy đi nào, Chúa tể, Kuon-senpai!!'' Nene hối thúc cả Kuon và Polka.

``Ờm, thôi. Mấy đứa cứ chạy đi,'' Kuon đáp, tay vẫn gõ lạch cạch.

``Sao anh bình chân như vại thế!?'' cô nàng ngỡ ngàng, rồi kéo tay Polka lại: ``Thôi, kệ anh! Ta chạy thôi, Chúa tể!''

``\eng{Tuân lệnh!!}'' nàng Cáo sa mạc gật đầu như tướng lĩnh nhận mệnh.

Họ lao ra ngoài như hai chiếc xe tăng mini, tay nắm chặt, chân đạp liên hồi. Từng bước gần hơn đến ánh sáng ở cuối đường hầm, từng giây tiến gần đến tự do\dots

%==========================================TRUYỆN PHỤ=========================================
\omk

Trong lúc Nene và Polka đã thoát khỏi cảnh tù đày\dots

Hai cô gái Thế hệ thứ Năm vừa trốn thoát khỏi cánh cửa lồng giam với tiếng thở dốc, toàn thân dính đầy bụi bặm. Polka thì vấp phải một cục gạch xốp gần cửa, còn Nene thì trượt chân do trời mưa. Nhưng cả hai không màng đến chuyện đó. Họ chạy như thể cuộc đời phụ thuộc vào nó, cuối cùng lăn quay ra thảm cỏ cạnh bãi xe của văn phòng.

\img{10-5f}

``\textbf{Chúng ta\dots\ sống rồi!}'' Nene hét lên như trong phim chiến tranh.

``Tự do\dots\ vị ngọt của tự do!'' Polka giơ hai tay lên trời, hưởng những giọt mưa rơi mát mẻ và ngọt lịm - một phần thưởng xứng đáng cho hai con người đã ``tái hòa nhập với xã hội.''

Một khoảnh khắc yên bình đến mức họ gần như quên rằng: người còn lại trong tù\dots\ vẫn chưa được cứu ra.

\begin{center}
  \uline{Văn phòng.}
\end{center}

``\dots Và đó là lúc bài hát của chúng ta kết thúc,'' Nene kết thúc phần giảng giải của mình tay cầm sáo.

Polka đứng ở bên cạnh khoanh tay, trước mặt cô là dàn mộc cầm: ``Mọi người hiểu chứ?''

Nàng Cáo sa mạc và nàng ``Hải cẩu'' vừa thoát khỏi ``nhà tù gỗ'' hồi sáng, giờ đã đứng trong văn phòng, bắt đầu trình bày kế hoạch âm nhạc ``có thể làm rúng động giới idol'' (theo lời Nene).

\img{10-5g}

Lamy (đang đứng) và Botan (ngồi bệt dưới sàn) đưa mắt nhìn nhau với ánh nhìn trống rỗng, mặt không biểu cảm như thể vừa bị ai đó bắt xem hết ba mùa ``P*ppa P*g'' mà không có phụ đề.

``Ở đoạn điệp khúc, chúng ta sẽ làm động tác `xé song sắt', sau đó là tư thế `bò ra cửa tù', rồi chuyển sang xoay tròn như đang tìm đường vượt ngục trong mê cung\dots'' Nene tiếp tục giảng giải.

Polka gõ nhẹ vào cây đàn xylophone, lớn giọng để gây chú ý: ``Còn phần solo của tôi là lúc giả tiếng còi hú báo động! Mọi người nhớ nốt này nè: ting-ting-ting-ting\dots\ \textbf{Hú hú!}''

Trong lúc hai cô nàng vẫn còn đang say sưa giải thích, hai cô nàng ở dưới hầu như không hề tỏ ra thích thú nào, nàng Sư tử trắng thậm chí còn ngáp ngắn ngáp dài.

``\hl{\dots Hình như họ định biến bản ballad cảm xúc thành nhạc nền phim hoạt hình sao, Shishiron?}''

``\hl{\dots Tôi cần reset não, Lamy ạ.}''

Sau một hồi trình bày như lên giáo án gdài như sớ Táo Quân, Nene cầm cây sáo recorder phát cho trẻ lớp 2 (có dán hình trái tim ở đầu), hô vang: ``\textbf{Vậy xin mời mọi người hãy thưởng thức!}''

Cô bắt đầu thổi\dots\ và đúng như bạn đoán: nốt nhạc bay tán loạn như kiến vỡ tổ, lên cao rồi rơi cái rụp, như thể sáo bị lỗi tông. Polka hòa theo bằng cách gõ loạn xạ lên xylophone, miệng hô: ``Câu này sẽ có trong bài kiểm tra đấy!''

Lamy rùng mình: ``Tôi tưởng bà nói bài hát về cảm xúc trong tù chứ\dots''

Botan thì giương mắt lờ đờ: ``Tôi tưởng là vượt ngục lãng mạn, không phải vỡ dạ dày âm thanh đâu\dots''

Cả hai đồng thanh cầu cứu trong tâm trí: ``\textit{Không biết Kuon-senpai đang làm gì nhỉ\dots/Kuon-senpai, làm ơn cứu bọn em với\dots}''

\begin{center}
  \uline{Bên trong nhà giam tạm thời.}
\end{center}

Kuon sau khi đã hoàn tất công việc trên laptop, khẽ duỗi vai, rồi ngửa người thở phào: ``\vie{Cuối cùng cũng xong việc. Không biết hai người họ diễn kịch xong chưa nhỉ\dots}''

Cậu quay sang bên cạnh. Trong chiếc lồng gỗ,  vốn là đạo cụ của một vở kịch cũ, giờ trống trơn, ngoại trừ vài thanh gỗ bị bẻ cong và một mảnh giấy nhỏ xíu có vẽ hình trái tim và dòng chữ nguệch ngoạc ``Polka ♥ Tự do''.

``\vie{Mấy đứa này\dots\ không biết là có cửa mở à?}'' cậu thở dài, lắc đầu cười mỉm.

Cậu ôm máy tính, tiến tới cửa gỗ lồng giam, rồi nhẹ nhàng vặn nắm cửa. Một tiếng *CẠCH* nhẹ. Cửa mở ngay tức khắc, như thể chưa từng bị khóa.

Bước ra ngoài, cậu thấy Watame vẫn đang say sưa ``thuyết giảng'', gương mặt sáng rỡ như thể đang lên giáo án cho môn Triết học Đại cương ba tiết liên tục. Còn Rushia thì\dots

Cô nàng xui xẻo lúc này đang ngồi bất động, đầu nghiêng một bên, mắt vô hồn. Trên má cô có viết nguệch ngoạc chữ ``HELP'' bằng son môi. Và miệng thì đang lẩm bẩm như tụng kinh:

\begin{center}
  ``Tù nơi đất lạ u sầu lắm,\\
  Chó cảnh bên tai nói\dots\ bậy bạ-nanodesu\dots''
\end{center}

``\dots Rushia-chan à\dots'' cậu cười méo xệch. ``Thơ của em có vần thật, mà nội dung thì thảm quá.''

Nàng Pháp sư cố quay đầu sang, ánh mắt tuyệt vọng, lẩm bẩm yếu ớt như một thây ma: ``Giúp em\dots\ Kuon-senpai\dots''

Cậu quay sang nàng Cừu vẫn đang luôn miệng nói, rồi thở dài một hơi. ``Ờm, Watame-chan, anh nghĩ là em nói đủ rồi đấy.''

Cậu tiến lại, vỗ nhẹ vào vai Cừu vàng. Vẫn không có gì cả.

``\dots và rồi nếu trộn nước cốt cam vào mochi matcha thì sẽ tạo ra vị mới\dots''

``Watame-chan? A-lô?''

``\dots nhưng nếu để lâu trong tủ lạnh thì lớp vỏ sẽ\dots''

``\dots Chán thật chứ.''

Cậu lấy từ túi áo ra một cái bịt mắt hình mèo, nhẹ nhàng đeo lên đầu nàng Cừu. Nhưng Watame vẫn nói tiếp như không có gì xảy ra. Dù mắt đã bị che, miệng vẫn hoạt động như máy in lỗi: liên tục và không kiểm soát.

Không còn cách nào khác, cậu quay sang Rushia, nói như ra lệnh: ``Xin lỗi Watame-chan nha, nhưng mà Rushia-chan sẽ đi cùng anh.''

Rushia gật đầu lia lịa như gà mổ thóc.

Cậu nhẹ nhàng kéo tay cô gái tóc xanh, cùng rời khỏi ``nhà tù kỳ quặc'' ấy, để lại ``Cerberus'' vẫn đang nói chuyện một mình giữa trời chiều thu se lạnh.

\ct

Ra đến hành lang, Kuon đưa cho Rushia một chai nước.

``Em ổn chứ?''

``Vâng\dots\ Em tưởng mình không bao giờ được nghe tiếng người nữa\dots''

``Rồi. Vậy cho anh hỏi, tại sao em lại bị bắt vậy?''

Rushia cúi đầu, rồi trả lời nhỏ như muỗi kêu: ``\dots Em ăn nhầm pudding trong tủ lạnh\dots''

``Của ai?''

``Marine-chan. Em mới chỉ ăn một miếng, vậy mà bị Fubuki-senpai bắt tại trận luôn-nanodesu.''

``\dots Đụng phải của Senchou thì xác định rồi. Anh cũng đã từng bị chính em ấy sỉ vả vào mặt vì tội `xúc phạm' món pudding\footnote{Xem lại {\flaregothic ÉPISODE 50}, quyển số Bốn.} đấy.''

``Em tưởng là hàng mẫu-nanodesu\dots''

Kuon thở dài: ``Đúng là cái nhà tù ngớ ngẩn thật\dots''

Nhưng cậu nào biết, toàn bộ sự việc - từ lồng giam, cai ngục, thuyết giảng và cả màn vượt ngục kỳ khôi - chỉ là một chiêu trò của Subaru, Marine và Fubuki để lấy ý tưởng cho bài hát mới.

Và khi bài hát ``Câu lạc bộ Ngục thất: Vượt ngục vì pudding'' được tung lên kênh \hll\ vào hai tuần sau, với beat trap nền nhạc lo-fi, lồng tiếng rên rỉ của Rushia và đoạn rap thuyết giảng của Watame\dots

Kuon chỉ biết gục mặt xuống bàn, lẩm bẩm: ``Tôi là nhân vật chính\dots bất đắc dĩ thật rồi.''

\ct

Cùng thời điểm mà Kuon đang dìu Rushia khỏi ngục tù, ở bên căn bếp, Marine đi tới tủ lạnh, hí hửng như trẻ con chuẩn bị mở quà Giáng Sinh. ``Hehe\~{} pudding hôm qua mua còn để dành. Nay trời lạnh lạnh ăn là đúng bài!''

Cô mở tủ lạnh và nhìn khay giữa.

\dots Trống trơn.

Nụ cười trên môi nàng Thuyền trưởng tắt ngúm, thay vào đó là một cái giật giật ở mắt phải. ``Hả\dots?''

Cô cúi xuống, nhìn các ngăn tủ dưới\dots\ Rrồi mở cả ngăn đá\dots

Không có một cái gì cả.

Marine gằn giọng: ``Nó\dots\ đâu rồi!?''

Cô vội chạy ra văn phòng, gặp A-chan đang cắm mặt vào máy tính. ``A-chan! Có ai lấy pudding của em không!? Em để tên rõ ràng bằng mực đỏ mà!!''

Nàng Đeo kính vẫn đang ngồi kiểm tra hòm thư, nhún vai đáp: ``Không rõ nữa. Thử hỏi bạn của em xem?''

Và thế là Cơn thịnh nộ Thuyền trưởng bắt đầu. Marine đi lùng mọi ngóc ngách, hỏi từng người. Khi nghe một người nói: ``À, chác là Rushia-chan ăn nhầm,'' thì\dots

``{\Huge \textbf{Rushiaaaaaaaaaa!!!}}''

Một tiếng hét bùng nổ hét lên giữa buổi chiều muộn, khiến lũ quạ bám trên cột điện giật mình bay mất.

Quả nhiên, vì quá cay cú, đích thân Marine đã thêm một đoạn ``cà khịa'' trong bài hát, trong đó có câu:

\begin{center}
  \uline{``Bà ăn pudding của tôi không xin phép,\\Giờ phải ở tù nghe Cerberus thuyết giảng suốt đêm\~{}''}
\end{center}

Tới độ khi mà cậu Tổng nghe bản demo của bài, cậu chỉ có thể rướn mày: ``Ê, ê\dots\ Anh vừa trở thành cameo trong một bài diss track à?''

\end{document}